\def\START{}\def\END{}
%% TEST-NAME: test-trinity-stripped
%% TEST-WHAT: Stripped integration fixture modelled on a real category-theory
%%            monograph (~16 000 lines). Exercises the full codependent.sty
%%            feature surface against the package stack found in real math
%%            papers: amsthm + thmtools (restatable), hyperref, cleveref,
%%            tikz-cd, custom math macros via \codepnewcommand, counter-
%%            sharing theorem environments, \codeptag, and concept tracking
%%            with def-site / use-site / forward-reference resolution.
%% TEST-SOURCE: Three-avatars monograph (Geiko–Mainiero–Moore);
%%              DESIGN.md §8a.0–§9a (integration surface);
%%              REVIEW_E all BLOCKERs composed
%% TEST-SECTION: DESIGN.md §8a-§9a (integration)
%% TEST-EXIT: 0
%% TEST-LOG-NOT: codependent.*Error|Undefined control sequence
%% TEST-LOG-NOT: Package codependent Warning
%% TEST-CDP-CONTAINS: \codep@cdp@end{OK}
%% TEST-CDP-CONTAINS: \codep@cdp@version{1}
%% TEST-CDP-CONTAINS: \codep@cdp@atom
%% TEST-CDP-CONTAINS: \codep@cdp@cmddef{Hom}{kind}{newcommand}
%% TEST-CDP-CONTAINS: \codep@cdp@cmddef{Hom}{arity}{2}
%% TEST-CDP-CONTAINS: \codep@cdp@cmddef{caty}{kind}{newcommand}
%% TEST-CDP-CONTAINS: \codep@cdp@cmddef{caty}{arity}{1}
%% TEST-CDP-CONTAINS: \codep@cdp@cmddef{Tr}{kind}{newcommand}
%% TEST-CDP-CONTAINS: \codep@cdp@cmddef{Tr}{arity}{0}
%% TEST-CDP-CONTAINS: \codep@cdp@use
%% TEST-CDP-CONTAINS: \codep@cdp@def
%% TEST-CDP-CONTAINS: \codep@cdp@label
%% TEST-CDP-CONTAINS: \codep@cdp@tag
%% TEST-CDP-COUNT: \codep@cdp@atom{1.2}{theorem} = 1
%% TEST-AUX-CONTAINS: \codep@atomref
%% TEST-AUX-CONTAINS: \codep@concept{Hom}
%% TEST-AUX-CONTAINS: \codep@concept{caty}
%% TEST-AUX-CONTAINS: \codep@concept{wtens}
%% TEST-AUX-CONTAINS: \codep@concept{Tr}
%% TEST-AUX-CONTAINS: \codep@concept{bimod}
%% TEST-AUX-CONTAINS: \codep@conceptref
%% TEST-ATOMS-MIN: 8
%% TEST-PACKAGES: amsmath,amsthm,thmtools,hyperref,cleveref,tikz,tikz-cd,codependent
%% TEST-RERUN: 2
%%
%% Modelled on a real operator-algebra / category-theory paper that uses:
%%   - thmtools \declaretheorem with counter sharing (sibling=)
%%   - thm-restate (\restatable + \restate)
%%   - cleveref (\cref, \Cref, multi-label)
%%   - hyperref (\autoref, \ref)
%%   - tikz-cd commutative diagrams (display + inline)
%%   - ~40 custom \newcommand macros for categories, functors, operators
%%
%% This fixture extracts the MINIMUM surface needed to exercise every
%% codependent feature against that stack. Content is paraphrased, not copied.
%%
%% Counter sharing: all theorem-like envs share a single counter named
%% "theorem" via thmtools sibling= key. This is the dominant pattern in
%% long math papers and is the primary stress case for atom numbering.

\documentclass{article}

%% ---------- Package stack (mirrors real paper) ----------
\usepackage[T1]{fontenc}
\usepackage{amsmath,amssymb}
\usepackage{mathtools}
\usepackage{amsthm}
\usepackage{thmtools}
\usepackage{thm-restate}

\usepackage{tikz}
\usetikzlibrary{cd}

\usepackage[hypertexnames=false]{hyperref}
\hypersetup{colorlinks=true,linkcolor=teal,citecolor=darkgray}
\usepackage{cleveref}

%% ---------- Theorem environments (counter-sharing) ----------
%% All theorem-like envs share the "theorem" counter, matching the
%% source paper's \makeTheoremEnv / \makeDefinitionEnv with sibling=.
\declaretheorem[
  name=Theorem,
  style=plain,
  numberwithin=section
]{theorem}
\declaretheorem[name=Proposition, style=plain, sibling=theorem]{proposition}
\declaretheorem[name=Lemma, style=plain, sibling=theorem]{lemma}
\declaretheorem[name=Corollary, style=plain, sibling=theorem]{corollary}
\declaretheorem[
  name=Definition,
  style=definition,
  sibling=theorem
]{definition}
\declaretheorem[name=Remark, style=remark, sibling=theorem]{remark}

%% ---------- codependent ----------
\usepackage[depth=3]{codependent}
\codeptrack{theorem,proposition,lemma,corollary,definition,remark}

%% ---------- Concept-tracked macros ----------
%% 5 macros converted to \codepnewcommand, modelled on the source
%% paper's operator-algebra notation. Each exercises a different arity.

% \Hom{X}{Y} — arity 2, the canonical concept-tracking example
\codepnewcommand{\Hom}[2]{\mathrm{Hom}(#1,#2)}

% \caty{C} — arity 1, category name formatter (mathsf in the source)
\codepnewcommand{\caty}[1]{\mathsf{#1}}

% \Tr — arity 0, trace operator
\codepnewcommand{\Tr}[0]{\operatorname{Tr}}

% \wtens{A} — arity 1, W*-bimodule tensor product
\codepnewcommand{\wtens}[1]{%
  \mathbin{\mathop{\otimes}\displaylimits^{\mathrm{w}}_{#1}}%
}

% \bimod{A}{M}{B} — arity 3, bimodule notation
\codepnewcommand{\bimod}[3]{{}_{#1}{#2}_{#3}}

%% ---------- Equation numbering by section ----------
\numberwithin{equation}{section}

\begin{document}
\START

%% ====================================================================
%%  Section 1: Overview (forward references to later definitions)
%% ====================================================================
\section{Three avatars of matrix product factors}%
\label{sec:avatars}

Given a pair of von Neumann algebras $A$ and $B$, one may consider
three closely related objects: matrix product factors, normal completely
positive maps, and pointed bimodules. Each is a ``face'' of the same
underlying categorical structure. Throughout, we write $\Hom{A}{B}$
for the morphism space in the ambient category $\caty{C}$.

\begin{theorem}\label{thm:equivalence}
  Let $A$ and $B$ be von Neumann algebras. The 2-categories
  $\caty*{MPF}^{\mathrm{op}}$ and $\caty{ptWBM}$ are locally adjoint
  equivalent.
\end{theorem}

\begin{proof}
  The identity-on-objects functor is constructed in
  \cref{def:bimodule_tensor}; see also \cref{lem:adjunction}.
\end{proof}

An intermediate paragraph after the theorem, referencing
\cref{thm:equivalence} and equation~\eqref{eq:tensor_compose}
below.

\begin{proposition}\label{prop:decomposition}
  Every normal completely positive map $\varphi\colon A \to B$
  admits a canonical decomposition into inner and outer parts.
  Specifically, if $\bimod{A}{M}{B}$ is the associated bimodule,
  we have
  \begin{equation}\label{eq:inner_outer}
    \varphi = \varphi_{\mathrm{in}} \circ \varphi_{\mathrm{out}}.
  \end{equation}
\end{proposition}

%% ====================================================================
%%  Section 2: Categorical framework (definitions + restatable)
%% ====================================================================
\section{Categorical framework}%
\label{sec:framework}

\begin{definition}\label{def:bimodule_tensor}
  Let $A$, $B$, and $C$ be von Neumann algebras. The
  $\Hom*{A}{B}$ set carries a natural bimodule structure. The
  \emph{bimodule tensor product} of $\bimod*{A}{M}{B}$ and
  $\bimod{B}{N}{C}$ over $B$ is the completion
  \begin{equation}\label{eq:tensor_compose}
    M \wtens*{B} N
    \;=\;
    \overline{M \otimes_B N}^{\,\|\cdot\|}.
  \end{equation}
  \codeptag{uid}{cat:bimodule-tensor}%
\end{definition}

A paragraph that references the definition: by
\cref{def:bimodule_tensor}, the tensor product is functorial.

\begin{lemma}\label{lem:adjunction}
  The functors $\mathtt{F}\colon \caty{MPF}^{\mathrm{op}} \to
  \caty{ptWBM}$ and $\mathtt{G}\colon \caty{ptWBM} \to
  \caty{MPF}^{\mathrm{op}}$ satisfy $\mathtt{F} \dashv \mathtt{G}$.
\end{lemma}

\begin{proof}
  The unit and counit are constructed from the canonical evaluation
  map; the triangle identities follow from $\Tr(\mathtt{id}) = 1$
  and the associativity of $\wtens{B}$.
\end{proof}

%% A tikz-cd display diagram: a commutative square of 2-functors
\begin{equation}\label{eq:square}
  \begin{tikzcd}
    \caty{MPF}^{\mathrm{op}}
      \arrow[r, "\mathtt{F}"]
      \arrow[d, "\mathtt{P}"']
    & \caty{ptWBM}
      \arrow[d, "\mathtt{K}"]
    \\
    \caty{CP}
      \arrow[r, "\mathtt{R}"']
    & \caty{WBM}
  \end{tikzcd}
\end{equation}

The square~\eqref{eq:square} commutes up to a natural isomorphism;
see \autoref{thm:equivalence}.

%% Inline tikz-cd diagram inside a paragraph
A short exact sequence
$\begin{tikzcd}[cramped, sep=small]
  0 \arrow[r] & K \arrow[r, hook] & M \arrow[r, two heads] & Q \arrow[r] & 0
\end{tikzcd}$
of bimodules witnesses the decomposition of
\cref{prop:decomposition}.

%% ====================================================================
%%  Section 3: Restatable theorem + restate
%% ====================================================================
\section{Morita product factors}%
\label{sec:morita}

Morita product factors generalise matrix product factors by replacing
finite-dimensional on-site spaces with arbitrary von Neumann algebra
modules.

\begin{restatable}{theorem}{MoritaEquiv}\label{thm:morita_equiv}
  The 2-category of Morita product factors is locally adjoint
  equivalent to the 2-category of pointed $W^*$-bimodules.
  Moreover, the equivalence restricts to the Type~I subcategory
  and recovers \cref{thm:equivalence}.
\end{restatable}

\begin{proof}
  The proof follows from a generalisation of \cref{lem:adjunction}
  to the Morita setting, combined with the canonical
  $\bimod{A}{L^2(A)}{A}$ bimodule.
\end{proof}

\begin{corollary}\label{cor:type_i}
  Every Type~I Morita product factor is a matrix product factor.
\end{corollary}

\begin{remark}\label{rem:trace}
  The trace $\Tr*$ plays a dual role: it normalises the inner product
  on $L^2(A)$ and defines the ``weight'' of a Morita product factor
  via $\Tr \circ \varphi$.
\end{remark}

%% ====================================================================
%%  Section 4: Restate appendix
%% ====================================================================
\section{Restated results}%
\label{sec:restate}

We collect restated results for the reader's convenience.

\MoritaEquiv*

The restatement above must not produce a duplicate sbl atom record.

A final paragraph referencing multiple labels via cleveref:
\Cref{thm:equivalence,prop:decomposition,lem:adjunction} together
establish the main correspondence, and
\cref{cor:type_i,rem:trace} describe the Type~I specialisation.

\END
\end{document}
